\section{Introduction}
An agent operating through a human-facing fixture can observe a button press,
contact event or haptic response. These observations are useful, but the inference
from ``contact happened'' to ``the authorized customer agreed'' is unjustified.
A controller can report a force pattern without identifying the toucher or
establishing willingness. A signed packet does not repair an untrusted sensor
path or make an old raw trace a fresh physical interaction.

OpenATT (Open Agent Tactile Trust) implements scoped controller evidence and
keeps it separate from permission. A stored challenge defines a short/long
contact pattern. A configured controller signs a bounded force/time trace and
the exact challenge. The verifier reprocesses the signal and checks binding,
current owner permission and required guarantees before creating one durable
local receipt. Contact reports never supply human consent or proximity.

This report provides concrete CSV acquisition, integer-unit signal processing,
controller signing, a shared transaction gate, attack cases and generated-trace
benchmarks. It does not report hardware sensing, human authentication, haptic
actuation or physical deployment. That distinction matters because the strongest
assumption is the trustworthy sensor-to-signing path, which the software profile
cannot create on its own.

\paragraph{Supported contribution.}
The artifact makes contact/consent separation and short-pattern limitations
executable. It supplies negative results: increased noise and timing variation
reduce acceptance, and an eight-bit pattern admits collisions across different
challenges. Signatures, hysteresis and database transactions are established
techniques. The equally expressive typed control matches the observed serial
policy behavior; no stronger authorization property or formal proof is claimed.
Research on sound composition and efficient evidence acquisition remains ahead.

